\section{Conclusiones y líneas futuras}
\subsection{Conclusiones}

El presente trabajo ha revisado el aprendizaje estructural de redes bayesianas basado en
restricciones, tomando el algoritmo PC como hilo conductor y articulando el análisis en
tres ejes: variantes del algoritmo, pruebas de independencia condicional y enfoques
híbridos. La conclusión principal es que el rendimiento depende de la interacción entre
potencia estadística, coste computacional y supuestos causales, y que no existe una
estrategia universalmente superior.

El PC es el método de referencia de la familia \emph{constraint-based}: cada decisión de
poda u orientación se apoya en una serie de contrastes de independencia condicional. Esa
trazabilidad es también su principal vulnerabilidad, pues en muestras finitas los errores
se propagan al esqueleto y a la orientación \citep{spirtes2000causation, colombo2014order}.
Las cinco limitaciones revisadas (consistencia en alta dimensión, dependencia del orden,
pérdida de potencia con $|\mathbf{S}|$, coste computacional y rigidez de supuestos) no se
eliminan con las variantes que tratamos, sino que se desplazan.

Las pruebas de independencia condicional determinan qué dependencias pueden detectarse y,
por tanto, qué grafos son alcanzables. La correlación parcial es eficiente bajo normalidad;
las tablas de contingencia son el contraste canónico para datos discretos, pero sufren
escasez de celdas; la regresión y la medida de covarianza generalizada cubren datos mixtos
y parte de la no linealidad, aunque dependen de la especificación de los modelos; y los
contrastes de kernel detectan dependencias no lineales a costa de un mayor coste y de una
calibración delicada \citep{shah2020gcm, zhang2011kci, runge2018conditional,
strobl2019approximate}. \citet{shah2020gcm} muestran que ningún contraste puede tener
potencia uniforme frente a todas las alternativas sin restringir la clase nula, lo que
confirma que no existe una prueba universal.

La independencia condicional tampoco identifica la dirección causal: bajo Markov y
fidelidad, el PC devuelve como máximo un CPDAG, por lo que parte de las orientaciones queda
indeterminada con datos observacionales \citep{spirtes2000causation}. Las extensiones
causales amplían el ámbito tratable, pero la orientación adicional exige supuestos de
separación, intervenciones o conocimiento de dominio
\citep{richardson2002ancestral, triantafillou2015constraint, mooij2020constraint}; la
inferencia causal es, por tanto, una inferencia bajo supuestos explícitos.

Los métodos híbridos combinan ambos paradigmas, pero no garantizan por sí solos un mejor
resultado. Las estrategias \emph{Restrict/Maximise}, desde CB y Sparse Candidate hasta MMHC
y H2PC, reducen el espacio de búsqueda, aunque un falso negativo en la fase de restricción
puede eliminar aristas irrecuperables
\citep{SINGH1993259, Friedman1999SparseCandidate, tsamardinos2006max, Gasse2014H2PC}.
ARGES muestra que la restricción debe ser adaptativa para preservar la consistencia
\citep{Nandy2018ARGES}, y GFCI incorpora variables latentes al combinar la búsqueda voraz
con las reglas de FCI \citep{Ogarrio2016GFCI}. El estudio comparativo de
\citet{Karadag2026Evaluation} muestra que ninguna familia domina sistemáticamente y que el
rendimiento depende de la topología, el tamaño muestral y la métrica. La clave no es
sustituir un paradigma por otro, sino diseñar la interfaz entre restricciones y
puntuaciones para que los errores de la primera fase sean recuperables por la segunda.

En definitiva, la pregunta relevante no es qué algoritmo es mejor en abstracto, sino qué
supuestos son defendibles en cada dominio y qué errores pueden tolerarse. La elección debe
guiarse por el tipo de dato, el tamaño muestral, la dimensionalidad, la densidad, la
presencia de variables latentes, el objetivo (predicción o estimación causal) y el
presupuesto computacional. El enfoque basado en restricciones mantiene una ventaja difícil
de sustituir: la posibilidad de auditar cada decisión y comunicar sus supuestos. Esa
interpretabilidad, junto con su papel en los métodos híbridos, explica por qué el PC y sus
variantes siguen siendo una referencia en el descubrimiento causal.

\subsection{Limitaciones del estudio}

Esta revisión es narrativa y no sistemática, por lo que la selección de trabajos puede estar
sesgada y las afirmaciones sobre rendimiento proceden de los estudios originales, no de una
reproducción independiente.

\subsection{Líneas futuras}

\begin{enumerate}
  \item \textbf{Garantías en muestra finita.} Desarrollar cotas, conjuntos de confianza
  sobre el CPDAG y procedimientos de calibración adaptativa del nivel de significación.
  \item \textbf{Pruebas adaptativas.} Seleccionar o combinar contrastes según el dato, el
  tamaño muestral y $|\mathbf{S}|$, con extensión a datos mixtos y de alta dimensión.
  \item \textbf{Supuestos relajados.} Integrar variables latentes, ciclos,
  no estacionariedad y eventos extremos con garantías formales, aprovechando datos
  intervencionales cuando estén disponibles.
  \item \textbf{Escalabilidad y enfoques diferenciables.} Combinar paralelización,
  algoritmos \emph{anytime} y relajaciones diferenciables que conserven garantías discretas
  e integren restricciones corregibles durante la optimización.
  \item \textbf{Evaluación reproducible.} Bancos de prueba comunes que varíen topología,
  muestra, dimensionalidad y confusión latente, informando de calidad estructural,
  calibración y coste computacional.
\end{enumerate}

Estas líneas son complementarias: las garantías condicionan el diseño de las pruebas, la
relajación de supuestos exige nuevas formas de evaluar la identificabilidad y la
escalabilidad determina qué métodos son aplicables. Su avance permitiría abordar el
descubrimiento causal como un problema de ingeniería de supuestos, datos y coste, en el que
el PC y sus variantes seguirían ocupando un lugar central.